\section{Камера, компьютерное зрение и рой}
\label{sec:python-vision}

\begin{infobox}[Цели главы]
\begin{itemize}
    \item Получать кадры с камеры дрона и находить на них маркеры ArUco и AprilTag.
    \item Вычислять положение дрона относительно маркера и строить на этом навигацию.
    \item Читать телеметрию датчиков в полёте и записывать её.
    \item Распознавать человека и жесты, управлять несколькими дронами одновременно.
\end{itemize}
\end{infobox}

\begin{warnbox}[Нужен настоящий дрон]
Симулятор отдаёт видео с виртуальной камеры и умеет кое-что из \texttt{cv2}: \texttt{detectMarkers} находит маркеры, расставленные на сцене, и возвращает их углы и номера. Но \texttt{cv2.solvePnP} в браузере всегда сообщает о неудаче, а библиотеки \texttt{mediapipe} нет совсем --- \texttt{import mediapipe} завершится ошибкой. Поэтому поиск маркеров можно попробовать в симуляторе, а оценку положения, трекинг человека и всё, что ниже, --- только на компьютере с реальным Пионером. Это ограничение среды, а не ошибка в программе.
\end{warnbox}
